\uuid{YTP1}
\exo7id{3949}
\titre{exo7 3949}
\auteur{quercia}
\organisation{exo7}
\datecreate{2010-03-11}
\isIndication{false}
\isCorrection{true}
\chapitre{Dérivabilité des fonctions réelles}
\sousChapitre{Autre}
\cpge{mpsi-mp2i.derivabilite.a, pcsi.derivabilite.a, ptsi.derivabilite.a, mpsi-mp2i.derivabilite.c, pcsi.derivabilite.c, ptsi.derivabilite.c, mpsi-mp2i.fonctions-deux-variables.b, pcsi.fonctions-deux-variables.b, ptsi.fonctions-deux-variables.b, bcpst1.derivation.a, bcpst1.derivation.b, bcpst1.fonctions-deux-variables.c}
\module{Analyse}
\niveau{L1}
\difficulte{}

\contenu{
\question{
Soit $f:\R \to \R$ dérivable telle que :
$\forall\ a,b \in \R,\ f(b) - f(a) = (b-a)f'\left(\frac {a+b}2\right)$.
Montrer que $f$ est polynomiale de degré inférieur ou égal à 2.
}
\reponse{
Dériver par rapport à $a$ puis par rapport à $b$.
}
}
